\subsection{Controlled representation readout (Stage-3 plan, Section~1.5)}
\label{sec:readout}
\textbf{Question and completed run.} The optional readout asks whether the causal
sites of L17H1 and L15H25 expose information about the mother's identity, and whether
that readout depends on the selected head. Run \texttt{stage3\_readout\_v1} contains
1{,}920 records: the 40 common pairs (20 discovery families, both orders), four
writer/site combinations, six conditions and two readout prompts. The frozen plan
matches the manifest hash; the gate passed with zero identity-injection and self-patch
drift and exact reproduction of three saved single-position importance probes. All
four analysis tables were regenerated locally and matched byte-for-byte. The 87
validation families were not used.

\textbf{Intervention and measurement.} Capture the residual at the output of the
writer's decoder block, including its MLP and shared normalization, and inject it at
the placeholder of a fixed readout prompt at the same block output. The prompts are
an identity completion (``cat $\to$ cat; \ldots; x $\to$'') and a
\texttt{mother\_of} question (``Who is the mother of this person?''). Conditions are
the intact clean source, replacement of the head at the capture site
(\texttt{head\_site}), replacement from the first differing token through that site
(\texttt{head\_span}), the intact corrupted source, the same token role in the other
answer-side fact (\texttt{irrelevant}), and no injection. The other-fact control
contains the competing mother's information; it is not an information-free vector.

Let $t_q,t_o$ be the first tokens of the query mother and the other candidate mother.
The primary statistic and paired change are
\[
 C=\log p(t_q)-\log p(t_o),\qquad
 \Delta C_c=C_c-C_{\mathrm{intact\ clean}}.
\]
A negative change moves the readout toward the swapped-in mother. Average the two
orders within each family, then average the 20 family values. The measurement concerns
first-token preference, not full-name generation. The source residual combines
multiple components; head dependence is assessed only through the paired intervention.

\begin{table}[htbp]
\centering\small
\setlength{\tabcolsep}{4pt}
\resizebox{\textwidth}{!}{%
\begin{tabular}{lllrrrr}
\toprule
Writer & Site & Readout & Clean $C$ & Source swap $\Delta C$ & Head-site $\Delta C$ & Negative families \\
\midrule
L17H1 & is & identity & $+0.51465$ & $-0.52070$ & $-0.45840$ & 20/20 \\
L17H1 & period & identity & $+0.28789$ & $-0.04609$ & $-0.04121$ & 20/20 \\
L17H1 & is & mother\_of & $+0.25459$ & $-0.00903$ & $-0.00981$ & 17/20 \\
L17H1 & period & mother\_of & $+0.26479$ & $-0.02686$ & $-0.02319$ & 20/20 \\
L15H25 & child last & identity & $+0.18213$ & $+0.02021$ & $+0.00605$ & 8/20 \\
L15H25 & child last & mother\_of & $+0.24312$ & $+0.00181$ & $-0.00020$ & 9/20 \\
L15H25 & is & identity & $+0.29395$ & $-0.15117$ & $+0.00303$ & 7/20 \\
L15H25 & is & mother\_of & $+0.25063$ & $-0.00171$ & $+0.00020$ & 5/20 \\
\bottomrule
\end{tabular}%
}
\caption{Controlled readout, family-level descriptives. Source swap uses the intact
corrupted prompt; head-site replaces only the selected head at the capture position.
The last column counts families with a negative head-site change after averaging
orders. The head-span condition is identical by construction at this capture layer
(see below). Values come from \texttt{readout\_family\_summary.csv}.}
\label{tab:readout}
\end{table}

\textbf{L17H1: identity-sensitive information, strongest at ``is''.} In the identity
readout, the contrast at ``is'' falls from $+0.515$ to $+0.056$ after head replacement,
compared with $-0.006$ for the corrupted source and $+0.002$ for the other-fact control.
The head-induced shift is negative in every family (median $-0.459$, range
$[-0.701,-0.234]$). Its magnitude is about 88\% of the mean full-source-swap shift.
This supports a contribution of the head's local output to the mother-token
preference exposed by this readout. The ratio compares two readout changes; it is
not a fraction of encoded information or of the original task's causal effect.

At the period, the identity shift is weaker but also negative in every family:
$-0.0412$ for head replacement versus $-0.0461$ for source swap, with median $-0.0410$
and range $[-0.0723,-0.0117]$. The \texttt{mother\_of} readout at the period also
responds consistently ($-0.0232$ versus $-0.0269$). These small responses do not by
themselves establish a relational binding. For comparison, the original-task
single-position effects are $+1.769$ logits at ``is'' and $+1.377$ at the period.
The head's approximately $5.7$-logit Scope-P effect is an all-position measurement,
not the effect at either site; none of these effects is numerically decomposed by
the readout shifts.

\textbf{L15H25: the child's-site content remains unresolved.} The child's last-token
site has original-task importance $+0.429$ logits, positive in every family. Replacing
the head changes the captured residual (mean norm difference 1.106), but the
\texttt{mother\_of} contrast barely changes: $-0.00020$ under head replacement,
and only $+0.00181$ under the full source swap. Identity readout also lacks a
consistent mother-swap response there. Since even the source-swap control does not
expose the expected information, the small head effect cannot distinguish missing
information from an insensitive readout. The query mother ranks first among the
candidates in 55\% of \texttt{mother\_of} probes for clean, head-replaced, corrupted
and no-injection conditions alike. Nor does identity readout reliably recover the
child: its first token ranks first in only 15\% of intact probes.

At L15H25's secondary ``is'' site, identity readout does respond to source swap
($-0.151$, negative in 18/20 families), while the head shift is approximately zero
($+0.003$), consistent with that site's small original-task effect ($+0.00394$).
This does not assign an address-only role to the head at its different, causal
child-last site.

\textbf{Readout quality and scope.} For every intact-clean probe, the top next token
is ``hello'' under identity and ``Mary'' under \texttt{mother\_of}; the experiment
has not freely decoded the names. At L17H1/``is'', identity readout assigns the query
mother's first token mean probability $0.000373$ (0.0373\%), compared with $0.0000429$
without injection, $0.000282$ for the other-fact residual, $0.000290$ after head
replacement and $0.000285$ for the corrupted source. Thus there is a positive
probability-relevance check and a source-specific contrast, but low absolute decoding
probability. These are means of per-example probabilities, not exponentiated mean
log probabilities. A null readout does not establish absent information; a
\texttt{mother\_of} prompt does not itself establish successful relation decoding.

\textbf{Why head-site and head-span are identical.} All 320 paired readout outputs
are exactly identical between these conditions. The head is replaced after its
layer's attention mixing, and the residual is captured at the output of that same
block. The remaining projection, normalization, residual addition and MLP act
independently at each position. Extra replacements at earlier positions therefore
cannot change the captured position within this block. This is a structural property
of the intervention, not evidence that earlier writes are unimportant. Testing their
propagation requires a later capture layer or a path intervention. The two conditions
must not be counted as independent corroboration.

\textbf{Implication.} The extension strengthens the hypothesis that L17H1 contributes
mother-identity-sensitive information at ``is'' and the period, while leaving the
representation at L15H25's child-last site unresolved. No writer-to-reader edge or
circuit has been demonstrated. Stage 4 should test consumers of these actual source
positions and retain L15H25 as a causal candidate despite the inconclusive readout.

